% 20-08(20쪽) — 두 함수 y=f(x)(왼쪽)·y=g(x)(오른쪽) 의 그래프. f 는 x<0 에서 y=1, f(0)=0, x>0 에서 y=-1.
% g 는 원점이 빈 점인 두 반직선(ㅅ 모양)이고 g(0)=-1. 스캔 화질이 낮아 TikZ 로 다시 그림.
\begin{tikzpicture}[line width=0.55pt, x=0.62cm, y=0.62cm, every node/.style={font=\footnotesize}]
  % 왼쪽 — y=f(x)
  \begin{scope}
    \draw[->, >=stealth, line width=0.7pt] (-2.2,0) -- (1.9,0) node[below=1pt] {$x$};
    \draw[->, >=stealth, line width=0.7pt] (0,-2.0) -- (0,1.9) node[left=1pt] {$y$};
    \draw[line width=0.6pt] (-1.9,1) -- (0,1);
    \draw[line width=0.6pt] (0,-1) -- (1.8,-1);
    \draw[fill=white, line width=0.5pt] (0,1) circle (2.2pt);
    \draw[fill=white, line width=0.5pt] (0,-1) circle (2.2pt);
    \fill (0,0) circle (2.2pt);
    \node[anchor=west, inner sep=2.5pt] at (0.14,1) {$1$};
    \node[anchor=east, inner sep=2.5pt] at (-0.14,-1) {$-1$};
    \node[below left=0pt and -2pt] at (0,0) {$\pt{O}$};
    \node[anchor=west] at (0.3,-1.7) {$y=f(x)$};
  \end{scope}
  % 오른쪽 — y=g(x)
  \begin{scope}[xshift=6.4cm]
    \draw[->, >=stealth, line width=0.7pt] (-2.1,0) -- (2.1,0) node[below=1pt] {$x$};
    \draw[->, >=stealth, line width=0.7pt] (0,-2.2) -- (0,1.3) node[left=1pt] {$y$};
    \draw[line width=0.6pt] (-1.8,-1.8) -- (0,0) -- (1.8,-1.8);
    \draw[fill=white, line width=0.5pt] (0,0) circle (2.2pt);
    \fill (0,-1) circle (2.2pt);
    \node[above left=-1pt and -2pt] at (0,0) {$\pt{O}$};
    % 원문대로 y축 왼쪽. inner sep 을 줄여 왼쪽 아래로 내려가는 반직선과 떨어뜨린다(2026-10-02 검수).
    \node[anchor=east, inner sep=1.2pt] at (-0.12,-1) {$-1$};
    \node[anchor=west] at (0.5,-2.0) {$y=g(x)$};
  \end{scope}
\end{tikzpicture}
